\documentclass[11pt,a4paper]{amsart}

\usepackage[T1]{fontenc}
\usepackage{amsmath,amssymb,amsthm,mathtools}
\usepackage{enumitem}
\usepackage{microtype}
\usepackage[margin=2.7cm]{geometry}
\usepackage{xcolor}
\definecolor{linkcol}{rgb}{0.10,0.20,0.50}
\usepackage[colorlinks=true,linkcolor=linkcol,citecolor=linkcol,urlcolor=linkcol]{hyperref}
\makeatletter
\g@addto@macro\UrlBreaks{\do\_\do\.}
\makeatother

\theoremstyle{plain}
\newtheorem{theorem}{Theorem}[section]
\newtheorem{proposition}[theorem]{Proposition}
\newtheorem{lemma}[theorem]{Lemma}
\newtheorem{corollary}[theorem]{Corollary}
\theoremstyle{definition}
\newtheorem{definition}[theorem]{Definition}
\newtheorem{example}[theorem]{Example}
\newtheorem{remark}[theorem]{Remark}

\newcommand{\CC}{\mathbb C}
\newcommand{\NN}{\mathbb N}
\newcommand{\Cu}{C_u^*}
\newcommand{\Ce}{C_e^*}
\newcommand{\Celoc}{C_{e,\mathrm{loc}}^*}
\newcommand{\OSadm}{\mathbf{OS}_{\mathrm{adm}}}
\newcommand{\OSex}{\mathbf{OS}_{\mathrm{ex}}}
\newcommand{\Esc}{\mathfrak E}
\newcommand{\Def}{\mathfrak D}
\newcommand{\Ideal}{\operatorname{Ideal}}
\newcommand{\Ch}{\operatorname{Ch}}
\newcommand{\ran}{\operatorname{ran}}
\newcommand{\id}{\operatorname{id}}
\newcommand{\coker}{\operatorname{coker}}
\newcommand{\Spec}{\operatorname{Spec}}
\emergencystretch=1.5em

\begin{document}

\title[Boundary migration and derived acyclicity]{Boundary migration, derived acyclicity, and countable continuity of $C^*$-envelopes}

\author{M.~Haufschild}
\thanks{Independent researcher. Email: \texttt{haufschildmichael@gmail.com}.}

\date{3 August 2026}

\subjclass[2020]{Primary 46L07, 46L30; Secondary 47L25, 46M18, 19K35}
\keywords{operator system, Shilov boundary, $C^*$-envelope, complete order quotient,
derived inverse limit, boundary representation, local $C^*$-algebra}

\begin{abstract}
Let $q:S\to T$ be a complete order quotient of separable unital operator systems, let
$\alpha:\Cu(S)\twoheadrightarrow\Cu(T)$ be the canonical map of universal $C^*$-covers,
and let $J_S,J_T$ be the Shilov ideals.  We introduce two canonical boundary-migration
algebras
\[
 \Esc(q)=\alpha(J_S)/(\alpha(J_S)\cap J_T),\qquad
 \Def(q)=J_T/(\alpha(J_S)\cap J_T).
\]
Escape vanishes exactly when $q$ descends to a surjection of $C^*$-envelopes; equivalently,
all target boundary representations pull back through the source envelope, or the corresponding
nc extreme boundary satisfies the spectral-closure criterion.  Under zero escape, deficit is the
exact failure of the envelope kernel to be generated by the operator-system kernel.  Its vanishing
is equivalent to positive matricial lifting along the Shilov ideals.

The principal result is a countable continuity law.  Zero-escape quotients form the maximal
category on which the $C^*$-envelope is covariantly functorial, and on this category the functor
preserves every countable inverse sequence after passage to Fr\'echet local $C^*$-envelopes:
\[
 \Celoc\!\left(\varprojlim S_n\right)\cong\varprojlim \Ce(S_n).
\]
At the level of universal Shilov ideals this is a derived acyclicity theorem.  We construct the
natural Roos exact sequence
\[
 0\longrightarrow (\varprojlim A_n)/(\varprojlim J_n)
 \longrightarrow\varprojlim(A_n/J_n)
 \longrightarrow\varprojlim{}^1J_n\longrightarrow0
\]
and prove that every Shilov system arising from a zero-escape tower satisfies
$\varprojlim{}^1J_n=0$.  This is not pointwise Mittag--Leffler exactness: for any prescribed
sequence $(D_n)$ of separable unital $C^*$-algebras we construct a tower of split zero-escape
complete order quotients with $\Def(q_n)\cong D_n$ at every link, while the global comparison
remains an isomorphism.  Deficits obey a canonical right-exact composition law.  We also obtain
faithful positive representations and a representable $K$-theory Milnor sequence for the local
envelope.
\end{abstract}

\maketitle

\section{Introduction}\label{sec:introduction}

The $C^*$-envelope of an operator system is its minimal generated $C^*$-cover
\cite{Hamana1979,Paulsen2002}.  It records the matrix order of the system while discarding the
largest boundary ideal; boundary representations give its concrete realization
\cite{DritschelMcCullough2005,Arveson2008}.  This minimality makes the construction rigid under
complete order embeddings, but quotient maps behave differently.  A
complete order quotient has positive lifts at every matrix level, up to the usual Archimedean
perturbation \cite{Han2011,KPTT2013}; nevertheless, it need not induce any map between the
corresponding $C^*$-envelopes.  The missing datum is not positivity but migration of the Shilov
boundary.

The first aim of this paper is to isolate that datum.  Let $q:S\to T$ be a complete order quotient
between separable unital operator systems.  Universality gives a surjective unital
$*$-homomorphism
\[
 \alpha:\Cu(S)\twoheadrightarrow\Cu(T)
\]
extending $q$.  If $J_S\triangleleft\Cu(S)$ and $J_T\triangleleft\Cu(T)$ are the Shilov ideals,
put
\begin{equation}\label{eq:defects-intro}
 I_q=\alpha(J_S)\cap J_T,\qquad
 \Esc(q)=\alpha(J_S)/I_q,\qquad
 \Def(q)=J_T/I_q.
\end{equation}
The algebra $\Esc(q)$ is source boundary mass that becomes visible in the target envelope.
The algebra $\Def(q)$ is target boundary mass that is not reached by source Shilov lifts.
They measure the two opposing inclusions between the source image and target Shilov ideals.

The escape algebra is the exact obstruction to functoriality.

\begin{theorem}[Boundary-migration criterion]\label{thm:intro-escape}
For a complete order quotient $q:S\to T$, the following conditions are equivalent.
\begin{enumerate}[label=\textup{(\roman*)}]
\item $\Esc(q)=0$.
\item $\alpha(J_S)\subseteq J_T$.
\item There is a surjective $*$-homomorphism
      $\overline q:\Ce(S)\twoheadrightarrow\Ce(T)$ extending $q$.
\item Every boundary representation of $T$, pulled back along $\alpha$, factors through
      $\Ce(S)$.
\item Under nc state-space duality, every nc extreme boundary point of the quotient state space
      belongs to the spectral closure of the nc extreme boundary of the source state space.
\end{enumerate}
\end{theorem}

Ordinary closure cannot replace spectral closure in the last condition.  The distinction is
essential beyond type~I covers; we use the factorization criterion of Kennedy--Kim--Manor
\cite[Proposition~9.2]{KKM2023}.

When escape vanishes, deficit becomes a cokernel in the category of $C^*$-algebras:
\[
 \Def(q)=J_T/\alpha(J_S).
\]
It has a concrete interpretation inside the source envelope.

\begin{theorem}[Kernel exactness and matricial lifting]\label{thm:intro-deficit}
Assume $\Esc(q)=0$.  Then
\begin{equation}\label{eq:deficit-kernel-intro}
 \Def(q)\cong
 \ker\overline q\big/\Ideal_{\Ce(S)}(\ker q).
\end{equation}
Moreover, the following are equivalent.
\begin{enumerate}[label=\textup{(\roman*)}]
\item $\Def(q)=0$.
\item $\alpha(J_S)=J_T$.
\item $\alpha|_{J_S}:J_S\twoheadrightarrow J_T$ is surjective.
\item For every $r\geq1$, every $b\in M_r(J_T)_+$ and every $\varepsilon>0$, there is
      $a\in M_r(J_S)_+$ such that
      $\alpha^{(r)}(a)=b$ and $\|a\|\leq\|b\|+\varepsilon$.
\end{enumerate}
\end{theorem}

The central theorem concerns inverse sequences.  Write $\OSadm$ for the category whose objects are
separable unital operator systems and whose morphisms are zero-escape complete order quotients.
Theorem~\ref{thm:intro-escape} makes $S\mapsto\Ce(S)$ a covariant functor on this category and
shows that no larger subcategory of complete order quotients admits the canonical envelope map.

\begin{theorem}[Countable continuity and Shilov acyclicity]\label{thm:intro-continuity}
Let
\[
 \cdots\longrightarrow S_3\longrightarrow S_2\longrightarrow S_1
\]
be a countable inverse sequence in $\OSadm$.  Put $A_n=\Cu(S_n)$, let $J_n\triangleleft A_n$ be
the Shilov ideal, and let $S=\varprojlim S_n$.  Then the canonical map is a local
$*$-isomorphism
\begin{equation}\label{eq:intro-continuity}
 \Celoc(S)\ \cong\ \varprojlim_n\Ce(S_n).
\end{equation}
Equivalently, the restricted inverse system of Shilov ideals is derived-acyclic:
\begin{equation}\label{eq:intro-acyclic}
 \varprojlim{}^1_nJ_n=0.
\end{equation}
The pulled-back stage boundary representations give a faithful local completely positive
representation of the algebra in \eqref{eq:intro-continuity}.
\end{theorem}

The vanishing in \eqref{eq:intro-acyclic} does not follow from surjectivity of the restricted maps
$J_{n+1}\to J_n$; these maps can fail as badly as possible.  The local envelope, rather than any
pointwise lifting property, forces global cancellation.

\begin{theorem}[Prescribed-deficit cancellation]\label{thm:intro-prescribed}
Let $(D_n)_{n\geq1}$ be any sequence of separable unital $C^*$-algebras.  There is an inverse
sequence in $\OSadm$
\[
 \cdots\twoheadrightarrow S_3\twoheadrightarrow S_2\twoheadrightarrow S_1
\]
such that every bonding map is split and
\[
 \Def(S_{n+1}\to S_n)\cong D_n
\]
for every $n$, while the global comparison \eqref{eq:intro-continuity} is an isomorphism.
\end{theorem}

Thus link deficits are local presentation data, not additive obstructions at infinity.  They do,
however, obey an exact finite-composition law.  If
$S\xrightarrow{q}T\xrightarrow{r}U$ are zero-escape morphisms with universal maps
$\alpha,\beta$, then
\begin{equation}\label{eq:intro-composition}
 \Def(q)\xrightarrow{\beta_*}\Def(rq)\longrightarrow\Def(r)\longrightarrow0
\end{equation}
is exact, and
\[
 \ker\beta_*=\frac{\alpha(J_S)+(J_T\cap\ker\beta)}{\alpha(J_S)}.
\]
This identifies the exact amount of an earlier deficit erased by the next quotient.

Theorem~\ref{thm:intro-continuity} has a useful homological shadow.  Phillips' representable
$K$-theory of $\sigma$-$C^*$-algebras \cite{Phillips1989} gives, for
$i\in\mathbb Z/2\mathbb Z$,
\begin{equation}\label{eq:intro-K}
 0\longrightarrow\varprojlim{}^1K_{i+1}(\Ce(S_n))
 \longrightarrow RK_i(\Celoc(S))
 \longrightarrow\varprojlim K_i(\Ce(S_n))\longrightarrow0.
\end{equation}
No extra term from boundary migration appears.

The proof has three layers.  Sections~\ref{sec:defects}--\ref{sec:composition} develop the
one-morphism defect calculus.  Section~\ref{sec:derived} proves a Roos exact sequence for an
arbitrary stable ideal tower.  Section~\ref{sec:continuity} uses local Shilov minimality to show
that the derived term vanishes for actual operator-system towers.  The construction in
Section~\ref{sec:prescribed} shows that this acyclicity is strictly global.  The last section records
faithful representations, $K$-theory, and the formal verification boundary.

The direction of the limit matters.  Mawhinney--Todorov developed inductive limits of operator
systems and obtained $C^*$-envelope results for particular graph systems
\cite{MawhinneyTodorov2018}.  Phillips developed inverse limits of $C^*$-algebras and their
$\sigma$-$C^*$ structure \cite{Phillips1988}, while Joi\c{t}a established the local Shilov theory
used here.  None of these results supplies an inverse-limit functoriality theorem for
$C^*$-envelopes, a derived Shilov sequence, or cancellation of arbitrarily prescribed link
deficits.  Inductive systems are governed by embeddings and Archimedeanization; the present
inverse systems are governed by quotient functoriality and derived failure of compatible Shilov
lifts.

\section{Preliminaries}\label{sec:preliminaries}

All operator systems are complex, unital, and norm closed.  A map is ucp if it is unital and
completely positive.  We use the operator-system quotient in the Archimedeanized sense of
\cite{KPTT2013}.  A surjective ucp map $q:S\to T$ is a \emph{complete order quotient} when the
induced map $S/\ker q\to T$ is a complete order isomorphism.  Equivalently, every positive matrix
over $T$ has the standard $\varepsilon$-positive lift; see \cite{Han2011}.

For an operator system $S$, its universal and minimal $C^*$-covers are denoted by $\Cu(S)$ and
$\Ce(S)$.  The quotient map
\[
 \epsilon_S:\Cu(S)\twoheadrightarrow\Ce(S)
\]
has kernel $J_S$, the Shilov ideal.  We identify $S$ with its canonical image in either cover.
For separable $S$, Arveson's boundary theorem gives
\begin{equation}\label{eq:boundary-kernel}
 J_S=\bigcap_{\pi\in\Ch(S)}\ker\pi,
\end{equation}
where $\Ch(S)$ denotes the boundary representations of $S$ in its universal cover
\cite{Arveson1969,Arveson2008}.

The quotient theorem of Kavruk--Paulsen--Todorov--Tomforde gives the universal-cover map used
throughout.

\begin{proposition}\label{prop:universal-quotient}
If $q:S\to T$ is a complete order quotient, there is a unique surjective unital
$*$-homomorphism $\alpha:\Cu(S)\to\Cu(T)$ extending $q$, and
\begin{equation}\label{eq:kernel-universal}
 \ker\alpha=\Ideal_{\Cu(S)}(\ker q).
\end{equation}
\end{proposition}

\begin{proof}
The ucp map $S\to T\subset\Cu(T)$ extends uniquely by universality.  Its image contains $T$ and
hence generates $\Cu(T)$, so it is onto.  Formula \eqref{eq:kernel-universal} is the universal
quotient identification obtained from \cite[Theorem~3.7 and Corollary~3.9]{KPTT2013}.
\end{proof}

We also use Fr\'echet local operator systems.  A countable inverse limit of $C^*$-algebras with
surjective bonding maps is a $\sigma$-$C^*$-algebra, hence a Fr\'echet locally $C^*$-algebra.
A local operator system is a unital self-adjoint subspace, equipped with the induced family of
matrix seminorms.  Local boundary ideals and local Shilov ideals are understood in the sense of
Joi\c{t}a \cite{Joita2025}.  For a separable Fr\'echet local operator system the local Shilov ideal
exists, is the intersection of local boundary-representation kernels, and quotienting by it gives
a generated local cover with trivial local Shilov ideal.  The latter cover is unique up to the
local $*$-isomorphism fixing the system \cite[Theorems~3.16 and~3.18]{Joita2025}.  The recent
$\Gamma$-boundary formulation of Joi\c{t}a--\c{S}imon gives an equivalent intrinsic kernel
description \cite{JoitaSimon2026}.

\section{Escape and boundary migration}\label{sec:defects}

Fix a complete order quotient $q:S\to T$ and its universal map $\alpha$.  Since the image of an
ideal under a surjective $*$-homomorphism is an ideal, both $\alpha(J_S)$ and $J_T$ are closed
ideals of $\Cu(T)$.  Definition \eqref{eq:defects-intro} therefore gives $C^*$-algebras.

\begin{remark}[The exchange algebra]\label{rem:exchange}
The relative position of the two ideals is packaged by
\[
 \mathfrak X(q)=\frac{\alpha(J_S)+J_T}{\alpha(J_S)\cap J_T}.
\]
The two summands have zero product modulo the intersection, and the addition map gives a canonical
$*$-isomorphism
\[
 \mathfrak X(q)\cong\Esc(q)\oplus\Def(q).
\]
Thus escape and deficit are the two orthogonal halves of one boundary-exchange algebra.
\end{remark}

\begin{proposition}\label{prop:escape-envelope}
The following are equivalent:
\begin{enumerate}[label=\textup{(\roman*)}]
\item $\Esc(q)=0$;
\item $\alpha(J_S)\subseteq J_T$;
\item $q$ extends to a surjective $*$-homomorphism
      $\overline q:\Ce(S)\to\Ce(T)$.
\end{enumerate}
When these conditions hold,
\begin{equation}\label{eq:envelope-map}
 \overline q(a+J_S)=\alpha(a)+J_T.
\end{equation}
\end{proposition}

\begin{proof}
The first two conditions are identical by definition.  The formula in
\eqref{eq:envelope-map} is well defined precisely when $\alpha(J_S)\subseteq J_T$.  It is then a
surjective $*$-homomorphism and extends $q$.  Conversely, any extension must satisfy
$\overline q\epsilon_S=\epsilon_T\alpha$ on the generating operator system and therefore on
$\Cu(S)$.  Hence $\alpha(J_S)\subseteq\ker\epsilon_T=J_T$.
\end{proof}

\begin{proposition}\label{prop:escape-boundary-reps}
Conditions in Proposition~\ref{prop:escape-envelope} are equivalent to: every
$\pi\in\Ch(T)$ satisfies $J_S\subseteq\ker(\pi\alpha)$.
\end{proposition}

\begin{proof}
By \eqref{eq:boundary-kernel},
\[
 \alpha(J_S)\subseteq J_T
 \quad\Longleftrightarrow\quad
 \pi(\alpha(J_S))=0\quad\text{for all }\pi\in\Ch(T).
\]
The latter condition says exactly that $\pi\alpha$ factors through $\Ce(S)$.
\end{proof}

\begin{proposition}\label{prop:escape-spectral}
Under the nc state-space duality, conditions in Proposition~\ref{prop:escape-envelope} are
equivalent to the spectral-closure condition in Theorem~\ref{thm:intro-escape}.
\end{proposition}

\begin{proof}
The quotient map identifies the nc state space of $T$ with the corresponding nc convex subset of
the nc state space of $S$.  Minimal $C^*$-covers are realized by restriction to the nc extreme
boundary.  By \cite[Proposition~9.2]{KKM2023}, an irreducible point representation factors through
the minimal cover exactly when its point lies in the spectral closure of that boundary.  Apply the
criterion to every nc extreme boundary point of the quotient state space and use
Proposition~\ref{prop:escape-boundary-reps}.
\end{proof}

Propositions~\ref{prop:escape-envelope}--\ref{prop:escape-spectral} prove
Theorem~\ref{thm:intro-escape}.

\begin{example}[A split quotient with escape]\label{ex:square-segment}
Let $K=[-1,1]^2$, $L=\{0\}\times[-1,1]$, and let $A(K)$ and $A(L)$ be the affine function
systems.  Restriction $q:A(K)\to A(L)$ is a split complete order quotient: the section extends an
affine function of $y$ constantly in $x$.  However,
\[
 \Ce(A(K))=C(\operatorname{ext}K)=C(\{(\pm1,\pm1)\}),\qquad
 \Ce(A(L))=C(\{(0,\pm1)\}).
\]
No $*$-homomorphism between these commutative envelopes can extend evaluation of affine
coordinates under restriction.  Thus $\Esc(q)\ne0$.  A ucp splitting alone does not control the
Shilov boundary.
\end{example}

\section{Deficit, matricial lifting, and realization}\label{sec:deficit}

Assume throughout this section that $\Esc(q)=0$.  Then
\begin{equation}\label{eq:deficit-zero-escape}
 \Def(q)=J_T/\alpha(J_S).
\end{equation}

\begin{proposition}[Algebraic deficit formula]\label{prop:deficit-formula}
There is a canonical $*$-isomorphism
\[
 \Def(q)\cong
 \ker\overline q\big/\Ideal_{\Ce(S)}(\ker q).
\]
\end{proposition}

\begin{proof}
Let $K=\ker\alpha$.  Proposition~\ref{prop:universal-quotient} gives
$K=\Ideal_{\Cu(S)}(\ker q)$.  Commutativity
\[
 \overline q\epsilon_S=\epsilon_T\alpha
\]
implies
\begin{equation}\label{eq:preimage-JT}
 \alpha^{-1}(J_T)=\epsilon_S^{-1}(\ker\overline q).
\end{equation}
Since $\alpha$ is onto, the third isomorphism theorem gives
\[
 \frac{J_T}{\alpha(J_S)}
 \cong\frac{\alpha^{-1}(J_T)}{J_S+K}.
\]
Quotienting \eqref{eq:preimage-JT} by $J_S$ identifies the numerator with
$\ker\overline q$, while
\[
 \epsilon_S(K)=\Ideal_{\Ce(S)}(\ker q).
\]
This yields the formula.
\end{proof}

\begin{proposition}[Shilov-matricial criterion]\label{prop:matricial}
Conditions in Theorem~\ref{thm:intro-deficit} are equivalent.
\end{proposition}

\begin{proof}
By \eqref{eq:deficit-zero-escape}, $\Def(q)=0$ exactly when
$\alpha(J_S)=J_T$, equivalently when $\alpha|_{J_S}$ is onto.  A surjective $*$-homomorphism of
$C^*$-algebras remains surjective on every matrix algebra and admits positive lifts with quotient
norm control.  This proves the lifting statement.

Conversely, positive lifting for $r=1$ implies surjectivity.  Every self-adjoint element is a
difference of two positive elements, and every element is a linear combination of two
self-adjoint elements.  Hence every element of $J_T$ is in the image of $J_S$.
\end{proof}

The next result shows that deficit has no hidden regularity constraint, even for split maps.

\begin{theorem}[Universal realization of deficits]\label{thm:universal-deficit}
Every separable unital $C^*$-algebra $D$ occurs as $\Def(q)$ for a split complete order quotient
$q:S\to\mathcal T_\infty$ with $\Esc(q)=0$, where
\[
 \mathcal T_\infty=\overline{\operatorname{span}}
 \{1,u_j,u_j^*:j\geq1\}\subset C^*(\mathbb F_\infty).
\]
\end{theorem}

\begin{proof}
The operator system $\mathcal T_\infty$ is the coproduct of countably many copies of
$\operatorname{span}\{1,z,z^*\}\subset C(\mathbb T)$.  It contains enough unitaries, and therefore
\begin{equation}\label{eq:Tinfty-envelope}
 \Ce(\mathcal T_\infty)=C^*(\mathbb F_\infty)
\end{equation}
by the enough-unitaries theorem \cite[Proposition~5.6]{Kavruk2014}; see also
\cite{FarenickEtAl2014}.  The countable statement follows from the compatible finite coproducts:
ucp maps are contractive and therefore extend to the norm closure.

Put
\[
 K=C([0,1],D),\qquad I=\{f\in K:f(0)=0\}.
\]
The ideal $I$ is essential.  Indeed, if $0\ne h\in H\triangleleft K$, continuity gives a point
$t>0$ with $h(t)\ne0$ (also when $h(0)\ne0$).  Choose $f\in C([0,1])$ with $f(0)=0$ and
$f(t)=1$.  Then $0\ne fh\in H\cap I$.

Choose contractions $k_j\in K$ with $k_1=0$ and
$C^*(1,k_2,k_3,\ldots)=K$.  The coproduct property gives a ucp map
\[
 \Phi:\mathcal T_\infty\to K,\qquad \Phi(u_j)=k_j.
\]
Let
\[
 B=C^*(\mathbb F_\infty)\oplus K,\qquad r(t)=(t,\Phi(t)).
\]
The range of $r$ generates $B$.  In fact,
\[
 r(u_1)^*r(u_1)-r(1)=(0,-1_K),
\]
so the generated algebra contains the central projection $(0,1_K)$.  It then contains
$(0,k_j)$ for every $j$, hence $0\oplus K$, and subtracting second coordinates gives
$C^*(\mathbb F_\infty)\oplus0$.

Define
\[
 S=r(\mathcal T_\infty)+(0\oplus I)\subset B,
 \qquad q(r(t)+(0,i))=t.
\]
The sum is closed.  If $x=r(t)+(0,i)$, the first coordinate gives
$\|t\|\leq\|x\|$, and then $\|i\|\leq2\|x\|$.  Projection $q$ and graph section $r$ are ucp, so
$q$ is a split complete order quotient with kernel $0\oplus I$.

The first projection is completely isometric on the graph and its quotient is
$C^*(\mathbb F_\infty)=\Ce(\mathcal T_\infty)$.  Hence the Shilov ideal of the graph copy is
exactly $0\oplus K$: any larger boundary ideal would descend to a nonzero boundary ideal in the
minimal cover.  If $H\triangleleft B$ is a boundary ideal for the larger system $S$, then
$H\subseteq0\oplus K$.  Complete isometry on the contained subspace $0\oplus I$ forces
$H\cap(0\oplus I)=0$.  Essentiality of $I$ gives $H=0$.  Hence
\[
 \Ce(S)=B.
\]
The induced envelope map is first projection.  Proposition~\ref{prop:deficit-formula} now gives
\[
 \Esc(q)=0,\qquad
 \Def(q)\cong K/I\cong D.
\]
\end{proof}

\begin{corollary}\label{cor:deficit-unbounded}
Split complete-order lifting imposes no dimensional, nuclearity, type~I, ideal-lattice,
representation-theoretic, or $K$-theoretic bound on deficit.
\end{corollary}

\section{The maximal functorial category and composition}\label{sec:composition}

\begin{definition}\label{def:admissible}
Let $\OSadm$ have separable unital operator systems as objects and zero-escape complete order
quotients as morphisms.  Let $\OSex$ be the wide subcategory of zero-deficit morphisms.
\end{definition}

\begin{theorem}\label{thm:maximal-category}
The category $\OSadm$ is the maximal subcategory of complete order quotients on which the
canonical $C^*$-envelope extension is defined.  The assignments
\[
 S\longmapsto\Ce(S),\qquad q\longmapsto\overline q
\]
form a covariant functor to unital $C^*$-algebras and surjective $*$-homomorphisms.
The zero-deficit morphisms form the wide subcategory $\OSex$.
\end{theorem}

\begin{proof}
Identity maps preserve Shilov ideals.  For composable maps with universal extensions
$\alpha,\beta$, zero escape gives
\[
 \alpha(J_S)\subseteq J_T,\qquad \beta(J_T)\subseteq J_U,
\]
and hence $(\beta\alpha)(J_S)\subseteq J_U$.  Universality identifies $\beta\alpha$ with the
universal map of the composite.  Formula \eqref{eq:envelope-map} then shows that envelope maps
compose.  Maximality is the converse implication in Proposition~\ref{prop:escape-envelope}: every
canonical envelope extension forces zero escape.

If both maps have zero deficit, then
\[
 (\beta\alpha)(J_S)=\beta(J_T)=J_U,
\]
so zero deficit is closed under composition.
\end{proof}

\begin{theorem}[Right-exact composition law]\label{thm:deficit-composition}
Let $S\xrightarrow{q}T\xrightarrow{r}U$ be morphisms in $\OSadm$.  There is a natural exact
sequence
\begin{equation}\label{eq:deficit-composition}
 \Def(q)\xrightarrow{\beta_*}\Def(rq)
 \longrightarrow\Def(r)\longrightarrow0,
\end{equation}
where
\[
 \beta_*(x+\alpha(J_S))=\beta(x)+\beta\alpha(J_S).
\]
Its first kernel is
\begin{equation}\label{eq:deficit-kernel-composition}
 \ker\beta_*=\frac{\alpha(J_S)+(J_T\cap\ker\beta)}{\alpha(J_S)}.
\end{equation}
Consequently \eqref{eq:deficit-composition} is short exact at the left after quotienting by
\eqref{eq:deficit-kernel-composition}, and it is already short exact when
$J_T\cap\ker\beta\subseteq\alpha(J_S)$.
\end{theorem}

\begin{proof}
Zero escape identifies the three deficits with
\[
 \Def(q)=J_T/\alpha(J_S),\quad
 \Def(rq)=J_U/\beta\alpha(J_S),\quad
 \Def(r)=J_U/\beta(J_T).
\]
The natural quotient from the middle algebra to the last is onto and has kernel
$\beta(J_T)/\beta\alpha(J_S)$, which is the image of $\beta_*$.  The kernel of $\beta_*$ consists
of those $x\in J_T$ for which $\beta(x)\in\beta\alpha(J_S)$, equivalently
$x\in\alpha(J_S)+(J_T\cap\ker\beta)$.  This proves both formulas.
\end{proof}

\section{The derived boundary sequence}\label{sec:derived}

We first separate the homological calculation from operator-system minimality.  Let
\[
 \cdots\xrightarrow{\alpha_2}A_2\xrightarrow{\alpha_1}A_1
\]
be a countable inverse sequence of complex vector spaces, with each $\alpha_n$ onto.  Let
$J_n\subseteq A_n$ be subspaces such that $\alpha_n(J_{n+1})\subseteq J_n$.  Put
$B_n=A_n/J_n$ and let $\overline\alpha_n:B_{n+1}\to B_n$ be the induced maps.

Define the Roos difference
\begin{equation}\label{eq:Roos}
 \Delta_X:\prod_{n\geq1}X_n\longrightarrow\prod_{n\geq1}X_n,
 \qquad
 \Delta_X((x_n))=(x_n-\alpha_n(x_{n+1}))_n
\end{equation}
for $X=A,J,B$, using restricted or induced maps as appropriate.  Then
\[
 \varprojlim X_n=\ker\Delta_X,\qquad
 \varprojlim{}^1X_n=\coker\Delta_X.
\]

\begin{theorem}[Derived boundary exact sequence]\label{thm:derived-exact}
There is a natural exact sequence
\begin{equation}\label{eq:derived-exact}
 0\longrightarrow
 \frac{\varprojlim A_n}{\varprojlim J_n}
 \xrightarrow{\Theta}
 \varprojlim(A_n/J_n)
 \xrightarrow{\partial}
 \varprojlim{}^1J_n
 \longrightarrow0.
\end{equation}
If the spaces are Fr\'echet, all maps are continuous and the last term carries the quotient
topology of $\coker\Delta_J$.
\end{theorem}

\begin{proof}
For $(\dot b_n)_n\in\varprojlim B_n$, choose representatives $b_n\in A_n$.  Compatibility gives
\begin{equation}\label{eq:mismatch}
 c_n=\alpha_n(b_{n+1})-b_n\in J_n.
\end{equation}
Replacing $b_n$ by $b_n+x_n$ with $x_n\in J_n$ changes $c$ to
$c-\Delta_Jx$.  Thus
\[
 \partial((\dot b_n)_n)=[c]\in\coker\Delta_J
\]
is well defined and linear.

The class $[c]$ vanishes exactly when $c=\Delta_Jx$ for some $x_n\in J_n$.  With
$a_n=b_n+x_n$, equation \eqref{eq:mismatch} becomes
$\alpha_n(a_{n+1})=a_n$.  Hence the compatible quotient family comes from
$\varprojlim A_n$.  Conversely, a coherent representative has zero mismatch.  This proves
exactness at the middle term.  The kernel of the map from $\varprojlim A_n$ is exactly
$\varprojlim J_n$, proving injectivity at the left.

To prove surjectivity of $\partial$, let $c=(c_n)\in\prod J_n$.  Choose $b_1=0$ and recursively
choose $b_{n+1}$ with
\[
 \alpha_n(b_{n+1})=b_n+c_n.
\]
This is possible because $\alpha_n$ is onto.  The quotient classes of $b_n$ are compatible and
have mismatch $c$.

For the topological assertion, let $C$ be the closed subspace of $\prod A_n$ consisting of
sequences coherent modulo $J_n$, and let $P=\prod J_n\subset C$.  The natural map
$C/P\to\varprojlim B_n$ is a continuous linear bijection between Fr\'echet spaces, hence a
homeomorphism by the open mapping theorem.  In this intrinsic model $\partial$ is induced by the
continuous mismatch map $C\to\prod J_n$, so it is continuous.
\end{proof}

\begin{remark}\label{rem:pointwise-not-global}
Surjectivity of every restricted map $J_{n+1}\to J_n$ makes $\Delta_J$ onto by recursive
correction and therefore kills $\varprojlim^1J_n$.  The converse is false even algebraically.  If
all restricted maps are zero, then every link has maximal cokernel while $\Delta_J$ is the
identity, so $\varprojlim^1J_n=0$.  The operator-system realization of this separation is much
stronger; see Theorem~\ref{thm:prescribed-tower}.
\end{remark}

\section{Countable continuity and Shilov acyclicity}\label{sec:continuity}

Let $(S_n,q_n)$ be an inverse sequence in $\OSadm$.  Write
\[
 A_n=\Cu(S_n),\qquad B_n=\Ce(S_n),\qquad J_n=\ker(A_n\to B_n).
\]
Let $\alpha_n:A_{n+1}\to A_n$ be the universal maps and
$\beta_n:B_{n+1}\to B_n$ the induced envelope maps.  Each map is onto.  Put
\[
 A=\varprojlim A_n,\qquad B=\varprojlim B_n,\qquad S=\varprojlim S_n.
\]

\begin{lemma}[Local Shilov ideal in the universal limit]\label{lem:local-shilov}
The local Shilov ideal of $S\subset A$ is
\begin{equation}\label{eq:local-shilov}
 J_{\mathrm{loc}}=
 \{a=(a_n)\in A:a_n\in J_n\text{ for every }n\}
 =\varprojlim J_n.
\end{equation}
\end{lemma}

\begin{proof}
Let $J^\infty$ denote the middle set in \eqref{eq:local-shilov}.  It is a closed ideal of $A$.
For $s\in M_r(S)$, the $n$th quotient seminorm modulo $J^\infty$ is bounded above by
$\|s_n\|$ and below by the quotient norm modulo the larger stage ideal $M_r(J_n)$.  Since $J_n$ is
a boundary ideal for $S_n$, the latter norm is $\|s_n\|$.  Thus $J^\infty$ is a local boundary
ideal.

Conversely, let $H\triangleleft A$ be any local boundary ideal.  The closed ideal
$\overline{p_n(H)}\triangleleft A_n$ is a boundary ideal for $S_n$: every matrix over $S_n$ lifts
to a matrix over $S$ because the bonding maps are onto, and the $n$th local quotient seminorm is
the quotient norm by $\overline{p_n(H)}$.  Hence $\overline{p_n(H)}\subseteq J_n$ for every $n$.
Therefore $H\subseteq J^\infty$, proving maximality and \eqref{eq:local-shilov}.
\end{proof}

Theorem~\ref{thm:derived-exact} now gives
\begin{equation}\label{eq:operator-derived}
 0\longrightarrow A/J_{\mathrm{loc}}
 \xrightarrow{\Theta}B
 \xrightarrow{\partial}\varprojlim{}^1J_n\longrightarrow0.
\end{equation}
The next lemma is the decisive boundary input.

\begin{lemma}[The inverse limit of stage envelopes is minimal]\label{lem:limit-minimal}
The local operator system $S$ locally generates $B$, and its local Shilov ideal in $B$ is zero.
\end{lemma}

\begin{proof}
Every coordinate projection $S\to S_n$ is onto.  Indeed, start with an element at stage $n$,
choose successive lifts under the surjective bonding maps, and define the lower coordinates by
composition.  Since $S_n$ generates $B_n$, a polynomial approximation to $b_n\in B_n$ lifts
generator by generator to a polynomial in elements of $S$.  Approximation in the $n$th norm
controls every earlier norm.  Hence $S$ locally generates $B$.

Let $H\triangleleft B$ be a local boundary ideal.  As in the proof of
Lemma~\ref{lem:local-shilov}, $H_n=\overline{p_n(H)}$ is a boundary ideal for
$S_n\subset B_n$.  But $B_n=\Ce(S_n)$ has zero Shilov ideal.  Thus $H_n=0$ for every $n$, and
therefore $H=0$.
\end{proof}

\begin{theorem}[Shilov derived acyclicity]\label{thm:shilov-acyclicity}
For every inverse sequence in $\OSadm$,
\[
 \varprojlim{}^1J_n=0,
\]
and the canonical comparison
\[
 \Theta:\Celoc(S)=A/J_{\mathrm{loc}}\longrightarrow B=\varprojlim\Ce(S_n)
\]
is a local $*$-isomorphism.
\end{theorem}

\begin{proof}
By Lemma~\ref{lem:limit-minimal}, $B$ is a generated local $C^*$-cover of $S$ with trivial local
Shilov ideal.  The quotient $A/J_{\mathrm{loc}}$ is another such cover.  Local envelope uniqueness
\cite[Theorem~3.18]{Joita2025} implements the identity on $S$ by a local $*$-isomorphism between
these covers.  Both maps agree on the generating local operator system with the coordinate
envelope quotients, so the isomorphism is precisely $\Theta$.  It is onto; exactness of
\eqref{eq:operator-derived} forces $\varprojlim^1J_n=0$.
\end{proof}

Theorem~\ref{thm:shilov-acyclicity} proves the algebraic and envelope assertions of
Theorem~\ref{thm:intro-continuity}.  Notice what it does not say: the maps
$J_{n+1}\to J_n$ need not be onto, and the coordinate images of $J_{\mathrm{loc}}$ need not equal
$J_n$.  A local isomorphism can recover a direction at a later seminorm without preserving the
chosen stage presentation term by term.

\section{Arbitrary deficits with global cancellation}\label{sec:prescribed}

We now prove the tower theorem promised in the introduction.

\begin{theorem}\label{thm:prescribed-tower}
For every sequence $(D_m)_{m\geq1}$ of separable unital $C^*$-algebras there is an inverse tower
of split morphisms in $\OSadm$
\[
 \cdots\twoheadrightarrow S_3\twoheadrightarrow S_2
 \twoheadrightarrow S_1=\mathcal T_\infty
\]
such that $\Def(S_{m+1}\to S_m)\cong D_m$ for every $m$.  Its local envelope is the inverse limit
of its stage envelopes.
\end{theorem}

\begin{proof}
Partition the free generators of $\mathcal T_\infty$ into an infinite sentinel set
$\{u_{s_m}:m\geq1\}$ and countably infinite blocks $G_m$.  Put
\[
 K_m=C([0,1],D_m),\qquad I_m=\{f\in K_m:f(0)=0\}.
\]
Each $I_m$ is essential and $K_m/I_m\cong D_m$.

Choose contractions $k_{m,j}\in K_m$ with the following properties:
\begin{enumerate}[label=\textup{(\alph*)}]
\item $k_{m,s_m}=0$;
\item $k_{m,s_\ell}=1$ whenever $\ell>m$;
\item $\{k_{m,j}:j\in G_m\}$ generates $K_m$ unitally.
\end{enumerate}
The coproduct property gives a ucp map
$\Phi_m:\mathcal T_\infty\to K_m$ with $\Phi_m(u_j)=k_{m,j}$.

We recursively construct $S_m\subset B_m=\Ce(S_m)$, a split quotient
$p_m:S_m\to\mathcal T_\infty$, and a ucp section
$r^{(m)}:\mathcal T_\infty\to S_m$.  Start with
\[
 S_1=\mathcal T_\infty,\qquad B_1=C^*(\mathbb F_\infty),\qquad
 p_1=r^{(1)}=\id.
\]
Given stage $m$, set $\Psi_m=\Phi_m p_m$ and
\begin{equation}\label{eq:tower-step}
 B_{m+1}=B_m\oplus K_m,\qquad
 S_{m+1}=\{(x,\Psi_m(x)):x\in S_m\}+(0\oplus I_m).
\end{equation}
The sum is closed by the first-coordinate estimate used in the proof of
Theorem~\ref{thm:universal-deficit}.  First projection
$q_m:S_{m+1}\to S_m$ is a split complete order quotient with graph section.  Put
$p_{m+1}=p_mq_m$ and let $r^{(m+1)}$ be the composite of $r^{(m)}$ with that section.

It remains to verify that $B_{m+1}$ is the stage envelope.  The sentinel
$v_m=r^{(m)}(u_{s_m})$ is a unitary in $B_m$: its base coordinate is $u_{s_m}$ and every earlier
sentinel coordinate is $1$ by condition~(b).  But $\Psi_m(v_m)=0$, so
\[
 (v_m,0)^*(v_m,0)-(1_{B_m},1_{K_m})=(0,-1_{K_m})
\]
belongs to the algebra generated by the graph.  This exposes the central $K_m$ summand.
Condition~(c) then gives $0\oplus K_m$, while the complementary central projection gives
$B_m\oplus0$.  Thus $S_{m+1}$ generates $B_{m+1}$.

The graph copy of $S_m$ has Shilov ideal $0\oplus K_m$.  A boundary ideal $H$ for the larger
system $S_{m+1}$ is therefore contained in $0\oplus K_m$.  Complete isometry on the contained
subspace $0\oplus I_m$ gives $H\cap(0\oplus I_m)=0$.  Essentiality of $I_m$ forces $H=0$, so
$\Ce(S_{m+1})=B_{m+1}$.

The envelope map of $q_m$ is first projection.  Its kernel is $0\oplus K_m$, while the ideal
generated by $\ker q_m$ is $0\oplus I_m$.  Proposition~\ref{prop:deficit-formula} gives
\[
 \Esc(q_m)=0,\qquad \Def(q_m)\cong K_m/I_m\cong D_m.
\]
Finally, Theorem~\ref{thm:shilov-acyclicity} applies to the resulting tower and proves the local
envelope assertion.
\end{proof}

\begin{remark}\label{rem:sentinel}
The future-sentinel condition is not cosmetic.  At step $m$, a generator used to expose the new
central projection must still be a full unitary in every previously added direct summand.  A
single repeatedly used generator would cease to be unitary after its first zero image.  The
countable reserve of sentinels closes this recursion without a hidden absorption assumption.
\end{remark}

\section{Representations, \texorpdfstring{$K$}{K}-theory, and formal verification}\label{sec:consequences}

\subsection{Faithful positive representations}

Let $(S_n)$ be an inverse sequence in $\OSadm$.  For each $n$, put
\[
 \rho_n=\bigoplus_{\pi\in\Ch(S_n)}\pi.
\]
Then $\ker\rho_n=J_n$ by \eqref{eq:boundary-kernel}.  If $\mathcal H_n$ is its Hilbert space, let
\[
 \mathcal H^{(N)}=\bigoplus_{n\leq N}\mathcal H_n,\qquad
 \mathcal D=\bigcup_N\mathcal H^{(N)}.
\]
Define the diagonal local representation by
\[
 \rho(a)|_{\mathcal H^{(N)}}=
 \bigoplus_{n\leq N}\rho_n(a_n).
\]
Every finite restriction is bounded and completely positive, and
\[
 \ker\rho=\bigcap_n p_n^{-1}(J_n)=J_{\mathrm{loc}}
\]
by Lemma~\ref{lem:local-shilov}.  Thus $\rho$ descends to a faithful local
$*$-representation of $\Celoc(S)$.  For one morphism, the same argument yields
\begin{equation}\label{eq:escape-detection}
 \Esc(q)\cong\rho_T(\alpha(J_S))\triangleleft\rho_T(\Cu(T))\cong\Ce(T),
\end{equation}
where $\rho_T$ is the direct sum of target boundary representations.  Escape is exactly source
Shilov mass detected by the faithful target boundary representation.

\subsection{Representable \texorpdfstring{$K$}{K}-theory}

Apply Phillips' Milnor sequence \cite[Theorem~3.2]{Phillips1989} to the $\sigma$-$C^*$-algebra
$\varprojlim\Ce(S_n)$ and use Theorem~\ref{thm:shilov-acyclicity}.  This proves
\eqref{eq:intro-K}.  The link deficits do not create a second derived term.  In the tower of
Theorem~\ref{thm:prescribed-tower}, they may nevertheless carry arbitrary prescribed $K$-theory.

\subsection{Formal verification boundary}

The exact quotient calculus has been formalized in Lean~4 against the repository's fixed mathlib
toolchain.  The formal module proves, without added axioms or placeholders:
\begin{itemize}
\item zero escape and zero deficit criteria for relative quotients;
\item exact finite matricial lifts and faithful-kernel detection;
\item the right-exact composition law and its kernel formula;
\item the Roos difference, the literal equivalence between compatible representatives and
      $\varprojlim(A_n/J_n)$, and the exact sequence \eqref{eq:derived-exact};
\item injectivity of the local comparison, surjectivity of the connecting map, and the equivalence
      between comparison surjectivity and vanishing of $\varprojlim^1J_n$;
\item recursive Mittag--Leffler correction and an explicit proof that linkwise cokernels need not
      control global $\varprojlim^1$.
\end{itemize}
The analytic operator-system layer uses the cited local Shilov existence and uniqueness theorems;
the formal module deliberately introduces no surrogate positivity or operator-system axiom.

\section{Concluding perspective}\label{sec:conclusion}

Complete order quotients possess all-level positive lifting, but that lifting does not determine
which boundary representations survive.  The exchange pair $(\Esc,\Def)$ separates two failures:
escape obstructs the existence of an envelope morphism, while deficit measures kernel
nonexactness once that morphism exists.  Escape therefore selects the maximal functorial category.

The countable law is more rigid than the one-link picture suggests.  On the maximal category,
$C^*$-envelopes commute with inverse limits and the universal Shilov ideal tower is always
$\varprojlim^1$-acyclic.  The prescribed-deficit theorem rules out a pointwise explanation: every
link may carry an arbitrary noncommutative obstruction, with arbitrary ideal theory and $K$-theory,
yet boundary minimality forces complete cancellation at infinity.  This is the universal boundary
migration law.

\begin{thebibliography}{99}

\bibitem{Arveson1969}
W.~Arveson,
\emph{Subalgebras of $C^*$-algebras},
Acta Math. \textbf{123} (1969), 141--224.

\bibitem{Arveson2008}
W.~Arveson,
\emph{The noncommutative Choquet boundary},
J. Amer. Math. Soc. \textbf{21} (2008), 1065--1084.

\bibitem{DritschelMcCullough2005}
M.~A.~Dritschel and S.~A.~McCullough,
\emph{Boundary representations for families of representations of operator algebras and spaces},
J. Operator Theory \textbf{53} (2005), 159--167.

\bibitem{FarenickEtAl2014}
D.~Farenick, A.~S.~Kavruk, V.~I.~Paulsen, and I.~G.~Todorov,
\emph{Operator systems from discrete groups},
Comm. Math. Phys. \textbf{329} (2014), 207--238.

\bibitem{Hamana1979}
M.~Hamana,
\emph{Injective envelopes of operator systems},
Publ. Res. Inst. Math. Sci. \textbf{15} (1979), 773--785.

\bibitem{Han2011}
K.~H.~Han,
\emph{On maximal tensor products and quotient maps of operator systems},
J. Math. Anal. Appl. \textbf{384} (2011), 375--386.

\bibitem{Joita2025}
M.~Joi\c{t}a,
\emph{The Shilov boundary for a local operator system},
Adv. Oper. Theory \textbf{10} (2025), no.~3,
\href{https://doi.org/10.1007/s43036-025-00459-5}{doi:10.1007/s43036-025-00459-5}.

\bibitem{JoitaSimon2026}
M.~Joi\c{t}a and G.-I.~\c{S}imon,
\emph{On the Shilov boundary ideal for Fr\'echet local operator systems},
arXiv:2605.02550 (2026).

\bibitem{Kavruk2014}
A.~S.~Kavruk,
\emph{Nuclearity related properties in operator systems},
J. Operator Theory \textbf{71} (2014), 95--156.

\bibitem{KPTT2013}
A.~S.~Kavruk, V.~I.~Paulsen, I.~G.~Todorov, and M.~Tomforde,
\emph{Quotients, exactness, and nuclearity in the operator system category},
Adv. Math. \textbf{235} (2013), 321--360.

\bibitem{KKM2023}
M.~Kennedy, S.-J.~Kim, and N.~Manor,
\emph{Nonunital operator systems and noncommutative convexity},
Int. Math. Res. Not. IMRN \textbf{2023} (2023), no.~5, 4408--4455.

\bibitem{MawhinneyTodorov2018}
L.~Mawhinney and I.~G.~Todorov,
\emph{Inductive limits in the operator system and related categories},
Dissertationes Math. \textbf{536} (2018), 1--57.

\bibitem{Paulsen2002}
V.~I.~Paulsen,
\emph{Completely bounded maps and operator algebras},
Cambridge Studies in Advanced Mathematics, vol.~78, Cambridge University Press, 2002.

\bibitem{Phillips1988}
N.~C.~Phillips,
\emph{Inverse limits of $C^*$-algebras},
J. Operator Theory \textbf{19} (1988), 159--195.

\bibitem{Phillips1989}
N.~C.~Phillips,
\emph{Representable $K$-theory for $\sigma$-$C^*$-algebras},
$K$-Theory \textbf{3} (1989), 441--478.

\end{thebibliography}

\end{document}
